\documentclass[11pt, a4paper]{article}
\input{preamble.tex}
\switchlang{ru}
\begin{document}

\title{1985 "--- Xerox Daybreak Mouse}
\date{}
\maketitle
\selectlanguage{russian}

В 1985 году компания Xerox выпустила официального преемника Xerox Star 8010 Information System (первого в мире офисного компьютера с оконным интерфейсом, появившегося в 1981 году) "--- рабочую станцию 6085 Professional Computer System, также известную как Daybreak. Компьютер поставлялся с версией 1.0 графической оболочки ViewPoint, основанной на программном обеспечении Xerox Star \cite{daybreak, hist}. Среди прочих улучшений, эти компьютеры получили и обновленную мышь.

\begin{figure}[h]
    \centering
    \includegraphics[scale=0.7]{1985_xerox_daybreak_mouse/pic_30.jpg}
    \caption{Xerox Daybreak Mouse}
    \label{fig:XeroxDaybreakPic}
\end{figure}

\section*{Внешний вид}

Мышь Xerox Daybreak не сильно отличается от оригинальной оптической мыши Xerox Star \cite{starmouse, photo}. Верхняя часть, выполненная из бежевого пластика, в определенной степени образована от параллелепипеда и по форме идентична мыши Xerox Star: она расширяется книзу, ее передняя и задняя стороны выпуклые и заметно наклонены, а вертикальные ребра закруглены, что существенно уменьшает ассоциацию с «коробкой» (рис. \ref{XeroxDaybreakTopAndBottom}). На верхней стороне корпуса расположены две большие продолговатые кнопки из более темного пластика.  Как и в случае с Xerox Star, они называются «SELECT» и «ADJUST» и, соответственно, используются для выделения и корректировки \cite{office}. 

\begin{figure}[h]
    \centering
    \includegraphics[scale=0.5]{1985_xerox_daybreak_mouse/top_30.jpg}
    \includegraphics[scale=0.5]{1985_xerox_daybreak_mouse/bottom_15.jpg}
    \caption{Xerox Daybreak Mouse, вид сверху и снизу}
    \label{XeroxDaybreakTopAndBottom}
\end{figure}

В отличие от мыши Xerox Star, расположение обеих кнопок смещено к краям корпуса, чтобы обеспечить возможность создания варианта мыши с третьей кнопкой между ними (для Xerox Star вариант мыши с тремя кнопками не предусматривался, чтобы максимально упростить управление, но в итоге третья кнопка оказалась полезной для инженерных и научных программных систем).

Нижняя часть мыши черная, на ней можно увидеть 3 пластиковые ножки, напечатанный номер модели мыши, а также выемку с линзами для светодиодов и матрицы оптических датчиков. Кабель мыши и корпус разъема выполнены в цвете кнопок. Для кабеля не предусмотрена механическая защита в месте выхода из корпуса мыши.

\begin{figure}[h]
    \centering
    \includegraphics[scale=0.3]{1985_xerox_daybreak_mouse/pad_30.jpg}
    \caption{Xerox Daybreak Mouse на реконструкции комплектного коврика}
    \label{fig:XeroxDaybreakPad}
\end{figure}

За регистрацию смещения цветовых неоднородностей при движении мыши отвечает сканирующая матрица размером $4 \times 4$ элемента, поэтому мыши требовался либо коврик со специально разработанным рисунком (массив светлых шестиугольников на темном поле), либо поверхность с похожим чередованием мелких светлых и темных пятен, например, джинсовая ткань. Коврик легко воспроизводим ксерокопированием, и более того, оригинальный коврик оптической мыши Star был бумажным и продавался в пачках по 25 листов \cite{pad}. Реконструированный коврик можно увидеть на рис. \ref{fig:XeroxDaybreakPad}.

\section*{Размеры}

По размеру мышь могла бы быть меньше и ниже, если бы не использовала стандартный корпус мышей Xerox Star, разработанный с учетом обеспечения одинакового внешнего вида для механических и оптических мышей. Тем не менее, по меркам 1980-х годов мышь Xerox Star можно назвать скорее компактной, чем средних размеров, а визуально она выглядит еще меньшей (рис. \ref{fig:XeroxDaybreakSize}).

\begin{figure}[h]
    \centering
    \includegraphics[scale=0.45]{1985_xerox_daybreak_mouse/size_15.jpg}
    \caption{Xerox Daybreak Mouse на размерном коврике с шагом сетки 1~см}
    \label{fig:XeroxDaybreakSize}
\end{figure}

\section*{Эргономика}

С точки зрения эргономики и внешнего вида, мышь выполнена в минималистичном стиле. Удобство использования мышей Xerox Star (включая мышь Daybreak), безусловно, меньше страдает от «прямоугольной» формы корпуса, чем это было с мышью Alto. Выпуклые продолговатые кнопки, расположенные в пределах досягаемости пальцев, по-прежнему удобны, а наклон ближней к пользователю стороны корпуса даже обеспечивает некоторый комфорт для ладони (рис. \ref{fig:XeroxDaybreakHand}) "--- по крайней мере, больше, чем у многих мышей 1980-х годов (конечно, более низкий корпус сделал бы мышь еще удобнее). Кроме того, небольшая ширина корпуса и его заметная высота, учитывая небольшой вес устройства, позволяет переносить мышь на другую часть коврика, удерживая ее с боков пальцами. Это могло бы быть полезным при кратковременном застревании механического блока из-за мусора "--- но в данном случае не особенно необходимо, учитывая, что мышь оптическая, а лист бумаги формата А4 представляет собой по-настоящему большой коврик.

\begin{figure}[h]
    \centering
    \includegraphics[scale=0.45]{1985_xerox_daybreak_mouse/hand_30.jpg}
    \caption{Xerox Daybreak Mouse с моделью руки человека}
    \label{fig:XeroxDaybreakHand}
\end{figure}

\section*{Подключение}

Как и все мыши Xerox, мышь Daybreak имеет квадратурный интерфейс \cite[p.~177]{aimouse}. В соответствии с исходной конструкцией Xerox Star, мышь подключается к компьютеру Daybreak через клавиатуру, которая имеет для этого специальный разъем на задней панели. Но, в отличие от предыдущих поколений компьютеров Xerox, вместо 15- или 9-контактного разъема D-SUB здесь используется фирменный миниатюрный штеккер нестандартного вида (рис. \ref{fig:XeroxDaybreakPic}). Однако гнездо на печатной плате клавиатуры показывает, что разница скорее косметическая, чем функциональная: технически это все еще 9-контактный разъем с шахматным порядком контактов, таким же, как у D-SUB \cite{jack, converter}.

\section*{Внутреннее устройство}

В разобранном виде манипулятор показан на рис. \ref{fig:XeroxDaybreakInside}, где можно увидеть специализированный чип в прозрачном пластиковом корпусе, изготовленном методом литья под давлением \cite{vlsi82}, который содержит матрицу сканирования, а также выполняет некоторую начальную обработку сигнала. Печатная плата поднята над основанием корпуса на вертикальных стойках, а под ней расположены под углом светодиоды для подсвечивания коврика мыши. Любопытно отметить, что кабель мыши Daybreak для обеспечения унификации с трехкнопочным аналогом включает провода для максимально возможного количества кнопок, причем провод третьей кнопки (зеленый на рис. \ref{fig:XeroxDaybreakInside}) даже аккуратно припаян к печатной плате.

\begin{figure}[h]
    \centering
    \includegraphics[scale=1]{1985_xerox_daybreak_mouse/inside_30.jpg}
    \caption{Xerox Daybreak Mouse в разобранном виде}
    \label{fig:XeroxDaybreakInside}
\end{figure}

\begin{thebibliography}{9}
\bibitem{daybreak} Xerox Daybreak: Wikipedia \url{https://en.wikipedia.org/wiki/Xerox_Daybreak}
\bibitem{hist} GUI Wonderland \#2: A Galaxy of Possibilities – Xerox Star \& Daybreak. - The Authopire Compendium, April 16th 2024 \url{https://blisscast.wordpress.com/2024/04/16/xerox-star-daybreak-gui-wonderland-2/}
\bibitem{office} Linden T., Harslem E. (editors). Office Systems Technology. A Look into the World of the Xerox 8000 Series Products: Workstations, Services, Ethernet, and Software Development, Xerox, November 1982. \url{https://www.bitsavers.org/pdf/xerox/sdd/OSD-R8203_Xerox_Office_Systems_Technology_Nov82.pdf}
\bibitem{starmouse} Degnan B. Xerox 8010 Star. - vintagecomputer.net 11/13/2011\url{https://vintagecomputer.net/browse_thread.cfm?id=311}
\bibitem {photo} Alves P. Xerox 6085 workstation from 1985. Retro Microcomputers, Workstations, Servers and Consoles - Welcome! November 27, 2018. \url{https://www.facebook.com/groups/retrocomputers/posts/2333452826684585/}
\bibitem{jack} Xerox 6085 photos. 1186 Keyboard 110K00890 -- bitsavers.org \url{https://www.bitsavers.org/pdf/xerox/6085/photos/1186_Keyboard_110K00890/1186_KB_4.JPG}
\bibitem{pad} R.\,F. Lyon The Optical Mouse: Early Biomimetic Embedded Vision / Advances in Embedded Computer Vision, Nov 2014, pp.3-22 \url{https://static.googleusercontent.com/media/research.google.com/ru//pubs/archive/43260.pdf}
\bibitem {vlsi81} R.\,F. Lyon. The Optical Mouse, and an Architectural Methodology for Smart Digital Sensors // VLSI DESIGN, August 1981. -- p. 20--30. \url{https://www.dicklyon.com/tech/OMouse/OpticalMouse-Lyon.pdf}
\bibitem {converter} Xerox Daybreak Mouse Converter \url{https://github.com/fiowro/XeroxDaybreakMouseConverter/}
\bibitem{aimouse} Xerox 6085 PCS Service Manual. October 1988 \url{https://www.bitsavers.org/pdf/xerox/6085/service/700P05710_6085_PCS_Service_Manual_Oct88.pdf}
\bibitem {vlsi82} R.\,F. Lyon, M.\,P. Haeberli. Designing and Testing The Optical Mouse // VLSI DESIGN, January/February, 1982. - p. 20--30. \url{https://www.dicklyon.com/tech/OMouse/DesigningTestingOMouse.pdf}
\end{thebibliography}
\end{document}
